\begin{abstractCN}
\setlength{\baselineskip}{23pt} %设置行距23磅

研究生学位论文是研究成果的最终呈现形式，其排版的规范与稳定，直接影响论文的可读性与评审体验。然而，深圳大学目前仅向研究生提供Word版本的模板，而Word在处理公式、长表格、跨章节交叉引用、参考文献编号等学术写作场景时，难以保证格式的一致性与可维护性。对于研究生而言，使用LaTeX完成论文已成为更为合理的选择。一个对齐学校最新规范、开箱即用的\LaTeX{}模板，因此具有显著的实用价值。

本模板的目标，是将《深圳大学研究生学位论文基本要求及格式规范》的核心条款封装于一个开箱即用的\LaTeX{}工程之中，使作者只需关注内容本身，而无须重复处理字体、间距、编号等格式细节。具体而言，本模板基于ctexart文档类，统一约束了页面边距、章节字体与字号、目录层级、公式与图表的逐节编号、脚注圈号样式，以及参考文献的引用与排版样式。同时，模板内置了中英文摘要、符号缩略语表、致谢、研究成果等环境，作者只需替换其中的占位文本，即可快速搭建出符合学校要求的论文框架。

为降低使用门槛，本模板采用“自文档化”设计：正文各章节本身即为模板的使用说明。作者编译生成的初始PDF，将依次呈现摘要写作要点、绪论中的图表与公式示范、各级标题的层级关系、参考文献的引用方式，以及致谢与研究成果的标准格式。在阅读PDF的过程中，作者即可同步理解每一处排版背后所对应的学校规范条款；着手写作时，只需将示范文字逐节替换为自己的研究内容，模板的格式约束将自动延续。

本模板由深圳大学计算机与软件学院的历届贡献者共同维护与迭代，并已在多届研究生的真实论文写作中得到验证。建议作者优先选择Overleaf在线平台进行编译，并将编译器设定为XeLaTeX；本地编译时，需确保系统已安装所需中西文字体。本模板以MIT协议开源，供深圳大学及兄弟院校自由使用，后续若学校规范更新，欢迎通过Issue或Pull Request协作维护。项目地址：\url{https://github.com/Cao-Wuhui/SZU_Latex_Master_Template}。

\end{abstractCN}


\begin{keywordCN}
LaTeX模板，学位论文，深圳大学，自文档化，XeLaTeX
\end{keywordCN}

\begin{abstractEN}
\setlength{\baselineskip}{23pt} %设置行距23磅

A graduate thesis represents the final form of one's research outcomes, and the consistency and stability of its typesetting directly affect both its readability and the reviewing experience. Shenzhen University, however, currently distributes only a Word template to its graduate students. Word struggles to maintain consistent formatting in academic scenarios such as mathematical equations, long tables, cross-chapter references, and bibliographic numbering, which makes it suboptimal for thesis writing. \LaTeX{} has therefore become the more reasonable choice for graduate students, and a ready-to-use \LaTeX{} template aligned with the university's latest specifications carries substantial practical value.

The goal of this template is to encapsulate the core clauses of the \textit{Basic Requirements and Format Specifications for Graduate Theses at Shenzhen University} into an out-of-the-box \LaTeX{} project, so that authors can focus on content rather than layout. Built upon the \texttt{ctexart} document class, the template consistently enforces page margins, sectional fonts and sizes, the hierarchy of headings, per-section numbering of equations, figures and tables, footnote styles, and the format of reference citations. Pre-built environments for the Chinese and English abstracts, the list of symbols and abbreviations, acknowledgments, and academic achievements are also provided, allowing authors to obtain a compliant thesis skeleton simply by replacing the placeholder text.

To lower the learning curve, this template adopts a self-documenting design: each chapter of the body text simultaneously serves as a usage guide. Compiling the initial PDF reveals, in sequence, the writing conventions for the abstract, demonstrations of figures, tables and equations in the introduction, the visual hierarchy of section levels, the citation conventions, and the standardized formats for the acknowledgments and academic achievements. Reading the PDF therefore provides both a working example and an explanation of the corresponding regulations; when the author begins writing, replacing the demonstrative text section by section is sufficient, and the formatting constraints will be inherited automatically.

This template has been collectively maintained by contributors from the College of Computer Science and Software Engineering at Shenzhen University, and has been validated across multiple cohorts of graduate theses. Authors are encouraged to compile the project on Overleaf with XeLaTeX as the engine; for local compilation, the required Chinese and Western fonts must be installed in advance. The template is released under the MIT License and is freely available to Shenzhen University and peer institutions. As the official specifications evolve, contributions via Issues or Pull Requests are warmly welcomed. Project repository: \url{https://github.com/Cao-Wuhui/SZU_Latex_Master_Template}.

\end{abstractEN}


\begin{keywordEN}
\LaTeX{} Template, Graduate Thesis, Shenzhen University, Self-Documenting, XeLaTeX
\end{keywordEN}

